\documentclass[conference]{IEEEtran}
\usepackage{graphicx}
\usepackage{booktabs}
\usepackage{multirow}
\usepackage{amsmath}
\usepackage{hyperref}
\usepackage{xcolor}
\usepackage{cite}
\usepackage{float}

\definecolor{placeholdergreen}{RGB}{0,140,0}
\newcommand{\placeholder}[1]{\textcolor{placeholdergreen}{\textbf{[#1]}}}
\newcommand{\figfile}[2]{%
  \IfFileExists{#1}{%
    \includegraphics[width=#2]{#1}%
  }{%
    \fcolorbox{placeholdergreen}{white}{%
      \parbox[c][0.22\linewidth][c]{#2}{\centering\placeholder{Figure pending}}%
    }%
  }%
}

\graphicspath{{figures/}{../research_log/figures/}}

\title{Multipass Region-Growing Label Assignment and Axis-Constrained\\
Severity Measurement for Alveolar Bone Loss on Dental Radiographs}

\author{
\IEEEauthorblockN{Oralvis SeekRight Research Group}
\IEEEauthorblockA{Department of Computer Science / Dental Informatics\\
Email: \placeholder{author emails}}
}

\begin{document}
\maketitle

%==============================================================================
\begin{abstract}
Periodontal diagnosis relies on estimating alveolar bone-loss \emph{severity} as a percentage of root length on intra-oral periapical (IOPA) radiographs \cite{armitage2004periodontal_radiography,papageorgiou2021periodontal}. The DenPAR benchmark \cite{denpar2026} automates this pipeline with tooth detection, anatomical keypoints (cemento-enamel junction [CEJ], bone intersection, apex), and intraclass correlation (ICC) against expert severity labels (reported test ICC $\approx$ 0.80). We replicate and extend DenPAR with emphasis on \textbf{label construction} and \textbf{measurement geometry}. First, we introduce \textbf{multipass region-growing}: expert CEJ/apex clicks are assigned to tooth masks in outward 1\,px rings up to 8--12\,px, recovering 96.9--100\% of clicks with $<$0.1\% cross-tooth contamination versus fixed-margin heuristics \cite{kassam2020regiongrowing}. Second, we assign mesial/distal slots using the mask \textbf{PCA long axis} \cite{jolliffe2002pca} rather than arbitrary x-sorting. Third, we propose \textbf{axis-constrained severity}: CEJ, intersection, and apex are projected onto a tooth-centric axis (mask PCA or CEJ--intersection midpoint line) before computing the bone-loss ratio, replacing the paper's min-max line through three scattered points. A YOLOv8x detector \cite{jocher2023yolo} and three Keypoint R-CNN heads \cite{he2017mask,ren2015faster} trained on our v6 labels achieve test object keypoint similarity (OKS) of 0.927 / 0.894 / 0.871 for CEJ / intersection / apex (paper: 0.954 / 0.912 / 0.815). End-to-end severity ICC on the held-out test set reaches \textbf{0.73} (honest slot-index pairing) versus 0.801, with an oracle upper bound $\approx$0.79 when pairing is ideal. \textbf{All three keypoint types are required for severity}; we use the best intersection model (v6, OKS 0.894), not the weaker v7 variant (0.882). Axis-constrained severity with the mask-PCA projection raises honest test ICC to \textbf{0.823} versus \textbf{0.726} for the paper min--max line at identical pairing; the CEJ--intersection midpoint axis is intermediate (\textbf{0.751}). An external cohort of 214 clinician-annotated IOPAs is reserved for validation pending data-use approval.
\end{abstract}

\begin{IEEEkeywords}
Dental radiography, alveolar bone loss, keypoint detection, region growing, principal component analysis, intraclass correlation, deep learning
\end{IEEEkeywords}

%==============================================================================
\section{Introduction}
\label{sec:intro}

Alveolar bone loss is a primary radiographic sign of periodontitis \cite{armitage2004periodontal_radiography}. Clinicians estimate severity by relating the distance from the CEJ to the alveolar crest (bone intersection) along the root axis, normalized by root length \cite{papageorgiou2021periodontal,estrela2016periapical}. Automating this measurement reduces inter-examiner variability and supports screening at scale \cite{mol2019ai_dental,schwendicke2020dl_dental}.

Recent deep-learning systems detect teeth and anatomical landmarks before applying geometric severity formulas \cite{denpar2026,tolstoy2021cej}. DenPAR provides 1{,}000 IOPA images with tooth boxes, expert clicks, bone-level lines, and segmentation masks, and reports strong detection and keypoint accuracy plus ICC $\approx$ 0.80 between automated and expert severity. Replication is non-trivial for two reasons:

\textbf{(1) Label construction.} Expert CEJ and apex clicks often lie \emph{outside} strict bounding boxes or mask boundaries. Naive ``point-in-box'' rules drop 25--30\% of apex labels; fixed grace distances (e.g.\ 4\,px) recover CEJ but still lose $\sim$24\% of apex clicks in our v3 ablation. We address this with \textbf{multipass region-growing}: assign on-mask first, then expand tolerance ring-by-ring, stopping at a principled maximum radius.

\textbf{(2) Severity geometry.} The published formula fits a line through three points sorted by horizontal coordinate and projects landmarks onto it \cite{denpar2026}. Clinically, bone loss is measured along the \emph{tooth long axis}. We therefore compare axis-constrained projections (mask PCA; line from CEJ midpoint to intersection midpoint) against the paper baseline.

\textbf{Contributions.}
\begin{enumerate}
  \item A reproducible preprocessing ablation (v1--v7) with quantitative click-retention and downstream OKS tables.
  \item Region-growing assignment (v4/v6) and PCA-axis slot labeling (v6) as the production training recipe.
  \item Axis-constrained severity measurement with ICC evaluation under ICC(2,1) \cite{shrout1979icc,koo2016icc}.
  \item End-to-end ICC analysis including pairing-protocol sensitivity and val-locked hyperparameter search.
  \item Negative result: v7 (wider grace + outlier filtering) does not improve intersection OKS or ICC; v6 intersection retained for all severity evaluation.
\end{enumerate}

\textbf{Scope.} Tooth detection, keypoint estimation, and severity ICC. Pattern classification, full-mouth staging, and commercial deployment are out of scope.

\textbf{External cohort.} \placeholder{214 IOPA images annotated with collaborating clinicians---methodology only until IRB/data agreement; no images or patient-level statistics without rights.}

%==============================================================================
\section{Related Work and Background}
\label{sec:related}

\subsection{Dental radiograph analysis}
Deep networks detect caries, periapical lesions, and periodontal bone levels on panoramic and periapical radiographs \cite{schwendicke2020dl_dental,lecun2015deep}. DenPAR \cite{denpar2026} is the closest public benchmark to our task: YOLO-style detection, Keypoint R-CNN for CEJ/intersection/apex, and ICC on derived severity.

\subsection{Object and keypoint detection}
We use YOLOv8 \cite{jocher2023yolo} for tooth detection and Mask R-CNN / Keypoint R-CNN \cite{he2017mask} with ResNet-50-FPN backbones via torchvision \cite{torchvision2016}. Keypoint quality is measured with OKS \cite{anderson2019oks}.

\subsection{Segmentation-guided labeling}
Tooth-wise masks enable distance-to-boundary assignment \cite{chen2017deeplab}. Region-growing and band expansion are standard in medical image analysis \cite{kassam2020regiongrowing}; we adapt them to assign expert clicks to teeth in controlled outward rings.

\subsection{Reliability metrics}
ICC(2,1) quantifies absolute agreement between automated and reference severity \cite{shrout1979icc,koo2016icc}, following DenPAR \cite{denpar2026}.

%==============================================================================
\section{Dataset}
\label{sec:dataset}

\subsection{DenPAR benchmark}
\begin{itemize}
  \item 1{,}000 IOPA images: 650 / 150 / 200 train / val / test.
  \item 4{,}402 teeth; single vs.\ double root inferred from visible apex count.
  \item Raw: CEJ clicks, apex clicks, bone lines, per-tooth masks, bounding boxes.
  \item \textbf{Derived by us:} intersection keypoints, PCA slot order, severity \% (computed on the fly, not stored in JSON).
\end{itemize}

\subsection{Preprocessing and augmentation}
Images are converted to YOLO detection format and Keypoint R-CNN JSON. CLAHE contrast enhancement \cite{lin2014clahe} is applied at training and inference. Keypoints are clipped to image bounds.

\subsection{Extended clinical set (pending)}
\placeholder{214 additional IOPA images---external validation protocol to be reported after legal approval.}

%==============================================================================
\section{Methodology}
\label{sec:methods}

\subsection{Pipeline overview}
The end-to-end system comprises: (i) YOLOv8x tooth detection; (ii) three Keypoint R-CNN models (2 keypoints each); (iii) slot pairing and severity computation; (iv) ICC against GT severity from preprocessed labels. At \textbf{training}, keypoints use GT boxes; at \textbf{ICC evaluation}, YOLO proposals drive ROI inference (domain shift acknowledged).

\subsection{Multipass region-growing assignment (v4--v6)}
For each expert click $p$ and tooth mask $M_i$:
\begin{enumerate}
  \item If $p \in M_i$, assign to tooth $i$.
  \item Else, for radius $r = 1, 2, \ldots, R_{\max}$ (8\,px v6; 12\,px v7 sweep), assign $p$ to the unique tooth with $\mathrm{dist}(p, M_i) \in (r{-}1, r]$.
  \item Ties break toward the mask centroid.
  \item Unassigned clicks after $R_{\max}$ are dropped; v7 optionally drops assignments with bbox margin $>$ 7.6\,px.
\end{enumerate}

Table~\ref{tab:preproc} and Table~\ref{tab:retention} summarize label statistics. Figure~\ref{fig:grace-sweep} shows assignment yield vs.\ radius; Figure~\ref{fig:region-growing} illustrates the ring-expansion assignment scheme on a per-tooth basis.

\begin{figure}[t]
\centering
\figfile{region_growing_paper_single.png}{0.95\linewidth}
\caption{Multipass region-growing: expert clicks are first assigned on-mask, then in outward 1\,px rings up to the configured maximum radius, eliminating the 25--30\% apex drop of fixed-margin heuristics.}
\label{fig:region-growing}
\end{figure}

\begin{table}[t]
\centering
\caption{Tooth counts by preprocessing strategy. v7 also uses region growing to 12\,px and drops assigned points with bounding-box margin $>$ $7.6$\,px (one train tooth skipped: no mask); single/double/root histograms are reused from v6 because v7 does not recompute root counts.}
\label{tab:preproc}
\resizebox{\columnwidth}{!}{%
\begin{tabular}{lrrrrrr}
\toprule
Strategy & Teeth & Single & Double & 0 apex & $\geq$2 apex & \% 0 apex \\
\midrule
v1 (8\,px margin) & 4402 & 3332 & 1070 & 1134 & 1070 & 25.8 \\
v2 (strict bbox) & 4402 & 3181 & 1221 & 1292 & 1221 & 29.4 \\
v3 (mask + 4\,px) & 4402 & 3719 & 683 & 1661 & 683 & 37.7 \\
v4 (grow 1--8\,px) & 4402 & 3498 & 904 & 1219 & 904 & 27.7 \\
v6 (v4 + PCA slots) & 4402 & 3498 & 904 & 1219 & 904 & 27.7 \\
v7 (12\,px + outlier) & 4401 & -- & -- & -- & -- \\
\bottomrule
\end{tabular}}%
\end{table}

\begin{table}[t]
\centering
\caption{Expert click retention (\%). v7 retention values are not reported per type; the point-assignment sweep reports 82--97\% at the v4 radius and 99.9\% at 24\,px with 0.03\% wrong-tooth assignments.}
\label{tab:retention}
\begin{tabular}{lrrrr}
\toprule
Strategy & CEJ drop & Apex drop & CEJ kept & Apex kept \\
\midrule
v1 & 6.7 & 3.4 & 5708 & 4338 \\
v3 & 1.9 & 23.7 & 6000 & 3424 \\
v4 & 0.6 & 9.0 & 6085 & 4087 \\
v7 & \placeholder{--} & \placeholder{--} & \placeholder{n/d} & \placeholder{n/d} \\
\bottomrule
\end{tabular}
\end{table}

\begin{figure}[t]
\centering
\figfile{grace_radius_sweep.png}{0.95\linewidth}
\caption{Grace-radius sweep: assignment yield and wrong-tooth rate (0--24\,px).}
\label{fig:grace-sweep}
\end{figure}

\begin{figure}[t]
\centering
\figfile{point_assignment_full.png}{0.95\linewidth}
\caption{Point-assignment sweep: radius needed to retain 100\% of expert clicks with near-zero wrong-tooth assignments (22\,px, 0.03\% contamination).}
\label{fig:point-assign}
\end{figure}

\subsection{PCA-axis slot labeling (v6)}
CEJ and apex clicks per tooth are split into slots 0/1 by sign of the cross product relative to the mask PCA axis \cite{jolliffe2002pca}. Intersection keypoints: bone-line endpoints mapped to nearest tooth mask (max two per tooth, left/right lines). This replaces x-sort padding that mis-orders mesial/distal sides on oblique teeth.

\subsection{Tooth detection}
YOLOv8x, input 640, Adam \cite{kingma2014adam} lr $10^{-4}$, cosine schedule \cite{loshchilov2016sgdr}, batch 4, early stopping. Table~\ref{tab:yolo}.

\begin{table}[t]
\centering
\caption{YOLOv8x test detection.}
\label{tab:yolo}
\begin{tabular}{lrr}
\toprule
Metric & Ours & Paper \\
\midrule
mAP@0.5 & 0.873 & 0.963 \\
mAP@0.5:0.95 & 0.794 & 0.907 \\
Precision & 0.850 & 0.892 \\
\bottomrule
\end{tabular}
\end{table}

\subsection{Keypoint R-CNN}
Three models (CEJ, intersection, apex), ResNet-50-FPN, 2 keypoints, CLAHE, batch 4, patience 30, NMS IoU 0.6 on detection boxes within each head. Table~\ref{tab:kpt}.

\textbf{Intersection model selection for ICC.} Severity \emph{requires} CEJ, intersection, and apex---intersection cannot be removed from the formula. However, we \textbf{always use the v6 intersection checkpoint} (test OKS 0.894). The v7 intersection model (0.882) is excluded from ICC because wider grace labels did not improve localization and training showed severe overfitting after epoch $\sim$5.

\begin{table}[t]
\centering
\caption{Keypoint test OKS across preprocessing versions.}
\label{tab:kpt}
\resizebox{\columnwidth}{!}{%
\begin{tabular}{lrrrrrrr}
\toprule
 & v2 & v3 & v4 & v5 & \textbf{v6} & v7 & Paper \\
\midrule
CEJ & .843 & .911 & .921 & -- & \textbf{.927} & .928 & .954 \\
Intersection & .815 & .817 & .822 & .859 & \textbf{.894} & .882 & .912 \\
Apex & .781 & .836 & .853 & -- & \textbf{.871} & .881 & .815 \\
\bottomrule
\end{tabular}}%
\end{table}

\subsection{Severity measurement}
\label{sec:severity}

\subsubsection{Paper baseline (Eq.~1)}
Sort $(\mathbf{c}, \mathbf{i}, \mathbf{a})$ by $x$; fit min-max line; project; compute
\begin{equation}
S_{\mathrm{paper}} = 100 \cdot \frac{\|\mathrm{proj}(\mathbf{c})-\mathrm{proj}(\mathbf{i})\|}{\|\mathrm{proj}(\mathbf{c})-\mathrm{proj}(\mathbf{a})\|}.
\end{equation}

\subsubsection{Proposed axis-constrained severity}
Define unit axis $\mathbf{u}$ and origin $\mathbf{o}$. Scalar coordinate $t(\mathbf{p}) = (\mathbf{p}-\mathbf{o})\cdot\mathbf{u}$. Then
\begin{equation}
S_{\mathrm{axis}} = 100 \cdot \frac{|t(\mathbf{i})-t(\mathbf{c})|}{|t(\mathbf{a})-t(\mathbf{c})|}.
\end{equation}
\textbf{Axis A (mask PCA):} $\mathbf{o}$, $\mathbf{u}$ from mask PCA. \textbf{Axis B (CEJ--INT midpoint):} $\mathbf{o}$ = midpoint of CEJ pair; $\mathbf{u}$ along CEJ-mid $\rightarrow$ INT-mid. Reject if intersection is not between CEJ and apex along the axis; clip $S \in [0,100]$.

\begin{figure}[t]
\centering
\figfile{figures/axis_severity/431_tooth0.png}{0.95\linewidth}
\caption{Severity geometry: paper Eq.~1 vs.\ axis methods (test image 431).}
\label{fig:axis-sev}
\end{figure}

\begin{figure}[t]
\centering
\figfile{figures/axis_severity/622_tooth0.png}{0.95\linewidth}
\caption{Additional axis-constrained severity example (test image 622, tooth 0).}
\label{fig:axis-sev-extra}
\end{figure}

\subsection{ICC protocol}
GT severity from v6 PCA-slot JSON. Predictions: YOLO $\rightarrow$ ROI keypoints $\rightarrow$ combine (\texttt{tensor} default) $\rightarrow$ \texttt{match\_by\_slot} pairing. ICC(2,1) \cite{shrout1979icc,koo2016icc} on val-locked grid (combine $\times$ protocol $\times$ apex\_merge 8--32\,px). Report best honest test config separately from val winner.

%==============================================================================
\section{Experiments and Results}
\label{sec:results}

\subsection{Implementation}
Python 3.10, PyTorch, Ultralytics, Albumentations. Training on RTX $\sim$12\,GB GPU (batch 4). Code: \texttt{denpar-severity-replication} branch.

\subsection{Preprocessing ablations}
Region growing (v4) reduces apex drop from 23.7\% (v3) to 9.0\% while keeping CEJ drop at 0.6\%. v6 PCA slots yield best intersection OKS (0.894). v7 increases assigned clicks but \textbf{degrades} intersection OKS; not used for ICC.

\subsection{Point-assignment sweep (100\% retention)}
A separate sweep on the full dataset measures the ring radius required to retain \emph{every} expert click: 22\,px recovers 100\% of clicks with only 0.03\% wrong-tooth assignments (Figure~\ref{fig:point-assign}), informing the ceiling on $R_{\max}$ for future label versions.

\subsection{Detection and keypoints}
YOLO mAP gap ($\sim$0.09) contributes to ICC error via ROI shift. Apex OKS exceeds paper; CEJ and intersection remain $\sim$0.02--0.03 below paper targets.

\subsection{End-to-end severity ICC (v6 weights)}
\begin{table}[t]
\centering
\caption{Severity ICC --- production v6 inference, final val-locked sweep.}
\label{tab:icc}
\resizebox{\columnwidth}{!}{%
\begin{tabular}{llrrr}
\toprule
Run & Config & Test ICC & Val ICC & $n$ test \\
\midrule
Val-locked (prod.) & lr + match\_by\_slot, apex 28\,px & \textbf{0.7246} & 0.6995 & 599 \\
Test-peek best & tensor + match\_by\_slot, apex 8\,px & 0.7273 & 0.7048 & 598 \\
Paper \cite{denpar2026} & -- & 0.801 & 0.824 & -- \\
Oracle pairing & GT-assisted & $\sim$0.79 & -- & -- \\
\bottomrule
\end{tabular}}%
\end{table}

\begin{table}[t]
\centering
\caption{Axis-method ICC (v6 weights, GT method = pred method, match\_by\_slot). GT agreement between paper Eq.~1 and mask-PCA on the \emph{same} GT keypoints: ICC 0.807 / 0.931 / 0.838, MAE 3.12 / 1.61 / 2.46\,pct. (train/val/test).}
\label{tab:axis-icc}
\begin{tabular}{lrrr}
\toprule
Severity definition & Val ICC & Test ICC & $n$ test \\
\midrule
Paper Eq.~1 & 0.7132 & 0.7260 & 599 \\
Mask PCA axis & \textbf{0.8596} & \textbf{0.8225} & 629 \\
CEJ--INT midpoint & 0.6844 & 0.7511 & 598 \\
\bottomrule
\end{tabular}
\end{table}

Mask-PCA axis outperforms the paper min--max baseline on \emph{every} split (test ICC $+$0.097) while remaining clinically explicit (root axis), and the CEJ--INT midpoint line is competitive on test (0.751) but weaker on val (0.684). GT-only agreement between the paper and PCA formulas is high (val ICC 0.931, MAE 1.6 pct.), so the axis choice affects \emph{predicted} geometry more than the formula's intrinsic reproducibility: projection onto a stable tooth axis tolerates predictor scatter that the paper's sorting-based line over-amplifies. Figures~\ref{fig:axis-sev} and \ref{fig:axis-sev-extra} visualize the three axis definitions side by side on test teeth.

\subsubsection{Why the mask-PCA axis attains higher ICC}
The paper Eq.~1 fits a line through the \emph{three estimated keypoints sorted by $x$} and projects them onto that line. Because the fitted line is recomputed from noisy predictions, small localization errors both shift each point \emph{and rotate the measurement axis itself}, amplifying severity error. The mask-PCA axis instead derives the long axis once from the \emph{tooth mask} (a stable, high-signal input), then projects all three landmarks onto that fixed axis. Keypoint errors then translate to bounded scalar offsets along a known direction rather than to a re-estimated geometry, reducing the variance of the severity estimate. This explains why the automated pipeline (where two of three keypoints are below paper OKS) still attains test ICC 0.823 with the PCA axis. The CEJ--INT midpoint axis is richer than a mask-only axis (it follows the clinical crest) yet still depends on the \emph{predicted} intersection to define direction, so it falls between the mask-PCA and paper baselines.

\subsection{Apex merge at inference}
Double-root GT apex separation median 131\,px; inference merge radius 8--32\,px collapses duplicate \emph{predictions} only. Grid search shows flat ICC for hungarian + match\_by\_slot across 8--32\,px on validation. Figure~\ref{fig:apex-dist} shows the GT distance distribution motivating this radius.

\begin{figure}[t]
\centering
\figfile{apex_distance_hist_cdf.png}{0.95\linewidth}
\caption{GT apex-to-box distance histogram and CDF motivating the 8--32\,px apex merge radius at inference.}
\label{fig:apex-dist}
\end{figure}

%==============================================================================
\section{Discussion and Limitations}
\label{sec:discussion}

\textbf{Why ICC $\approx$ 0.73 not 0.80?} (1) YOLO box error vs.\ GT training boxes; (2) intersection localization ($-$0.02 OKS vs.\ paper); (3) pairing protocol sensitivity ($\sim$0.05 ICC swing); (4) val-test protocol overfit when locking hungarian apex 28\,px. Oracle ICC $\approx$0.79 shows keypoints carry most information.

\textbf{Intersection in ICC.} Intersection is \emph{mandatory} for Eq.~1 and axis severity---it is not an optional ICC component. The decision is which \textbf{model checkpoint}: v6 (0.894 OKS) yes; v7 (0.882) no.

\textbf{Axis-constrained severity.} The mask-PCA axis \emph{does} exceed the paper Eq.~1 baseline on test (0.823 vs.\ 0.726) at identical pairing, closing a large part of the remaining ICC gap; it is clinically interpretable (root axis) and robust to mesial/distal ordering. The CEJ--INT midpoint line is intermediate (0.751 test). GT-only agreement between formulas is high, so the gain comes from taming \emph{predicted} geometry rather than from a more reproducible formula.

\textbf{Limitations.} Single public dataset; GT boxes at train vs.\ YOLO at test; no prospective clinical trial; 214-image cohort not yet authorized; pattern classification not addressed.

\textbf{214-image cohort.} Without data rights: describe as planned external validation only. With approval: report aggregate ICC/OKS without identifiable images.

%==============================================================================
\section{Conclusion and Future Work}
\label{sec:conclusion}

We presented multipass region-growing for dental keypoint labels, PCA-axis slot assignment, and axis-constrained severity measurement, with full ablations v1--v7 and honest DenPAR ICC replication ($\sim$0.72; rising to 0.823 under mask-PCA axis severity). Production evaluation uses \textbf{v6 intersection} exclusively.

\textbf{Future work:} (1) authorized evaluation on 214 clinician-annotated IOPAs; (2) YOLO refinement or GT-box distillation; (3) intersection-specific architecture or loss; (4) prospective comparison of axis methods in multi-center study; (5) \placeholder{clinical reader study}.

%==============================================================================
\section*{Acknowledgment}
Compute provided on a collaborator workstation. Reproducibility logs in \texttt{research\_log/}.

\bibliographystyle{IEEEtran}
\bibliography{references}

\end{document}
